\section{Reinforcement learning for investment strategy}\label{sec4}

{\color{blue}This section describes the decision layer of MGAT-FinRL and its interface with the forecaster. The aim is to measure how much the forecasts change the allocations learned by standard DRL agents when everything else in the environment is held fixed.\par}

\subsection{FinRL-based portfolio optimization}\label{FinRL-based portfolio optimization}

{\color{blue}As illustrated in Fig.~\ref{fig:finrl}, MGAT is trained first and then held fixed; its out-of-sample forecasts are part of the agent's state at every decision in the training, selection, and evaluation periods.\par}
\begin{figure*}[!t]
    \centering
    \includegraphics[width=\textwidth]{fig2_finrl.pdf}
    \caption{{\color{blue}FinRL layer of MGAT-FinRL. The state combines the 252-day covariance matrix and four technical indicators computed from observed prices with two forecast rows produced by MGAT at the close of the decision day. The agent's action is projected onto long-only weights, which are executed at the next close in a share-and-cash ledger with trading costs. The lower bar shows the walk-forward use of the data.}}
    \label{fig:finrl}
\end{figure*}

We evaluate the decision-making layer under the FinRL framework with a compact set of widely used continuous-control algorithms: PPO\citep{schulman2017ppo}, A2C\citep{mnih2016a3c}, SAC\citep{haarnoja2018sac}, TD3\citep{fujimoto2018td3}, and DDPG\citep{lillicrap2016ddpg}. They span the main design axes in practical portfolio control. PPO and A2C are on-policy actor-critic methods with stable updates under non-stationary return distributions; PPO further constrains policy updates by clipping and can use entropy regularisation to encourage exploration. DDPG and TD3 are off-policy deterministic actor-critic methods for continuous control; TD3 mitigates Q-value overestimation through target policy smoothing and delayed actor updates. SAC is an off-policy stochastic method that maximises expected return plus an entropy term, which promotes exploration. All agents use the Stable-Baselines3 implementations with the FinRL portfolio-allocation interface \citep{Liu2021FinRL}; the core algorithms are not modified.

\paragraph{Coupling with MGAT}
{\color{blue}At the close of day $t$, MGAT forecasts $\hat q_{t,i}$ (Eq.~\eqref{eq:q}) and the high and low of day $t+1$ for every stock. Two rows are added to the state: the cross-sectional $z$-score of the predicted tradable-horizon return,
\begin{equation}
z_{t,i}=\frac{\hat q_{t,i}-\frac1N\sum_{j}\hat q_{t,j}}{\mathrm{sd}_j(\hat q_{t,j})},
\label{eq:zq}
\end{equation}
and the predicted log range $\hat r_{t,i}=\hat y^{h}_{t,i,1}-\hat y^{l}_{t,i,1}$, standardised with its mean and standard deviation over the RL training period. Each forecast is out of sample: it is produced by the forecaster of the walk-forward fold whose test year contains day $t+1$ (Section~\ref{subsec:splits}). Because the forecaster is held fixed, portfolio rewards do not update MGAT, and all agents and seeds receive identical forecasts.

The same environment is run with forecasts from four sources, each producing the same two rows: MGAT; the graph-free Transformer trained with the same objective; the ridge regressions of Section~\ref{subsec:objective}; and a naive source with $\hat q_{t,i}=0$ (so $z_{t,i}=0$) and the observed range $\log(H_{t,i}/L_{t,i})$ of day $t$ in place of the predicted range. The naive source fixes the state dimension and removes forecast information, so differences between the MGAT and naive agents measure the contribution of the forecasts to the learned allocations.\par}

\paragraph{State space}
{\color{blue}The state at decision time $t$ is
\begin{equation}
s_t=\left[\boldsymbol{\Sigma}_{t};\ \mathbf{D}_{t};\ \mathbf{z}_{t}^{\top};\ \tilde{\mathbf{r}}_{t}^{\top}\right]\in\mathbb{R}^{(N+6)\times N},
\label{eq:state}
\end{equation}
where $\boldsymbol{\Sigma}_{t}\in\mathbb{R}^{N\times N}$ is the sample covariance matrix of the last 252 daily close-to-close returns, $\mathbf{D}_{t}\in\mathbb{R}^{4\times N}$ stacks MACD, the 30-day RSI, the 30-day CCI, and the 30-day DX of each asset computed with the \texttt{stockstats} library from observed daily high, low, and close prices up to day $t$, $\mathbf{z}_t$ is given by Eq.~\eqref{eq:zq}, and $\tilde{\mathbf{r}}_t$ is the standardised predicted range. For $N=15$, the state has $21\times15$ entries for every forecast source.\par}

\paragraph{Action space and feasibility constraints}
At each decision time $t$, the policy outputs a real-valued vector $\mathbf{c}_{t}\in[-1,1]^{N}$, which is mapped to a long-only portfolio weight vector using a softmax projection:
\begin{equation}
w_{t,i}=\frac{\exp(c_{t,i})}{\sum_{j=1}^{N}\exp(c_{t,j})},\qquad
\mathbf{w}_{t}=(w_{t,1},\ldots,w_{t,N})\in\Delta^{N},
\label{eq:portfolio_softmax}
\end{equation}
where $\Delta^{N}=\{\mathbf{w}\in\mathbb{R}^{N}\mid w_i\ge 0,\ \sum_{i=1}^{N}w_i=1\}$. {\color{blue}With actions bounded in $[-1,1]$, the largest weight ratio between two stocks is $e^{2}\approx7.4$ and a single weight lies between 0.010 and 0.35, so the policies tilt an approximately diversified portfolio rather than concentrate it.} All agents operate on the same continuous action space.

\paragraph{Reward and training objective}
{\color{blue}
At decision time $t$ the agent chooses target weights $\mathbf{w}_t$ using information available at the close of day $t$, and the trade is executed at the close of day $t+1$ (Section~\ref{subsec:env}). Let $V_t$ denote the account value at the close of day $t$, after any trade executed on that day. The immediate reward is the scaled net one-step portfolio return,
\begin{equation}
R_{t+1}=\kappa_R\,\frac{V_{t+1}-V_t}{V_t},\qquad \kappa_R=0.1,
\label{eq:rl_reward}
\end{equation}
where $V_{t+1}$ is computed from observed prices and is net of the trading costs charged on day $t+1$ (Eq.~\eqref{eq:trading_cost}). The agents maximise the expected discounted sum of rewards with discount factor $\gamma=0.99$. Each training episode covers the RL training period once, starting from an all-cash account of $V_0=10^6$.

A common alternative in FinRL-style environments is to return the updated account value, $R_{t+1}=V_{t+1}$, as the immediate reward. Wealth-level rewards are non-stationary in scale and weight late, high-wealth steps more heavily, whereas the one-step return in Eq.~\eqref{eq:rl_reward} has a stable scale across the episode and lets trading costs enter the learning signal directly. Reward variants that penalise volatility or drawdown were not evaluated.
\par}
